\section{Related Work}
\label{sec:related}

A finite two-tier pool poses two ordering problems: eviction ranks who
leaves, and placement ranks who sits in the fast tier. Both oracles
reduce to B\'el\'ady next-reference distance, which is not observable
online, so every policy substitutes a proxy conditioned on progressively
richer information.

\subsection{Information-Conditioned Eviction Methods}
\label{sec:related-evict}

Access-history methods treat a recent touch as evidence of future need.
PagedAttention~\cite{kwon_pagedattention_2023} pages a growing KV pool
in the style of virtual memory, and
SGLang~\cite{zheng_sglang_2024} reuses shared prefixes through a radix
tree. Both evict by block recency, so a session alive on a tool wait
appears oldest and leaves first.

Lifecycle methods promote the pause to a first-class request state.
InferCept~\cite{abhyankar_infercept_2024} chooses discard, preserve, or
swap from a per-request waste account, and
Continuum~\cite{li_continuum_2026} pins paused KV with a time-to-live
derived from reload cost. Both keep a paused session resident yet issue
an absolute keep-or-drop for one request, so the locally cheap choice
can be the globally expensive victim.

Arrival-table methods predict the next request.
AGSERVE~\cite{agserve_2025} reads the next arrival from a tool-type
latency table, assigning the same wait to every session that shares a
tool, which homogenizes heterogeneous jobs and mis-ranks the longest
residents.

Structure-conditioned methods exploit workflow or role diversity.
CacheScout~\cite{cachescout_2026} learns a Markov table over agent
roles, KVFlow~\cite{pan_kvflow_2025} reads steps-to-execution off a
declared graph, and PBKV~\cite{zheng_pbkv_2026} predicts future agent
calls. The contrast in all three is borrowed from structural diversity;
a homogeneous batch flattens the scores and the policy falls back to
recency.

A parallel line shortens the token set of one trajectory rather than
ranking who leaves a shared pool.
IntentKV~\cite{li_intentkv_2026} scores history tokens from cross-turn
intent, MemDecay~\cite{matam_memdecay_2026} assigns region-specific
decay, and CommitKV~\cite{huang_commitkv_2026} retires pages only after
a tool-call commit. All three trade accuracy for a smaller working set
without comparing live sessions against one another.
Semantic queues~\cite{fang_saecache_2026}, variable-grain
paging~\cite{jeon_granikv_2026}, and workflow-aware job
admission~\cite{ni_topas_2026} similarly reshape what a request
contains or which job runs, yet none ranks live sessions on a shared
pool.

Learned replacement policies trained by reinforcement learning or
imitation, such as LRB~\cite{song_lrb_2020} in web caching, achieve
strong offline accuracy but require per-decision inference times of
hundreds of microseconds to milliseconds, which exceeds the
single-digit microsecond budget of a near-memory scheduler by two
orders of magnitude and therefore cannot serve as an online eviction
oracle in the hardware pipeline targeted here.

Taken together, access history inverts the ranking, a per-request
lifecycle yields no ranking, a population table homogenizes it, and
borrowed structure yields a ranking that vanishes when the batch is
uniform, and none of these proxies conditions on the signal that every
session carries in every regime, namely its own observed gaps and its own
progress toward completion.

\subsection{Hierarchy-Aware Placement Strategies}
\label{sec:related-place}

A kept session still has an access cost; the oracle places the soonest
next reference in the fast tier. Prior placement work has widened its
information from a static offload plan to job-scheduler hints to a
global reuse store.
FlexGen~\cite{sheng_flexgen_2023} solves a linear program over GPU,
CPU, and disk, conditioning on membership in the current computation, so
a session alive on a tool wait is demoted first.
CachedAttention~\cite{gao_cachedattention_2024} places blocks from
scheduler hints across a three-level store, but the hint names the next
scheduled job, not the remaining wait of each session.
Mooncake~\cite{qin_mooncake_2025} pools DRAM, SSD, and RDMA under
prefix-popularity ordering, which ranks by reuse frequency rather than
by proximity.

Consequently, a static plan inverts the placement in time, a scheduler
hint homogenizes it across pauses, and a popularity score orders it by
frequency instead of proximity, so all three keep the right bytes in
the wrong place because none conditions on the gap this session just
opened.

\subsection{Near-Memory and Datapath KV Hardware}
\label{sec:related-hw}

Token, device, and interconnect engines decide which bytes of one
request to compute or move, a decision complementary to ranking who
stays in a shared pool. Table~\ref{tab:hw-kv} in
Section~\ref{sec:hweval} places both objects on one grid.

Importance datapaths prune or skip tokens inside one decode step.
Token-Picker~\cite{park_tokenpicker_2024} withholds transfers based on
pre-softmax probability, UniCAIM~\cite{xu_unicaim_2025} prunes via
content-addressable search, a KV-MMU~\cite{moradifirouzabadi_kvmmu_2025}
replaces least-relevant tokens at runtime,
HiKV~\cite{fang_hikv_2026} evicts tokens then loads significant
elements through a reconfigurable sorter, and
Kelle~\cite{xia_kelle_2025} evicts by attention score in embedded DRAM.
Look-back and voting designs such as MATA~\cite{zhu_mata_2026} and
VEDA~\cite{wang_veda_2025} belong to the same datapath family.
Capacity-expansion hardware enlarges the store.
CXL-SpecKV~\cite{liu_cxlspeckv_2026} offloads KV to FPGA memory
with speculative prefetch, and V-Rex~\cite{kim_vrex_2026} retrieves
clustered video-token subsets. None of these engines observes a tool
wait or a pool of other live sessions.

A near-memory scheduler sits beside the hierarchy and ranks session
residency on request and gap events, so that eviction and tiering
become two reads of that ranking. The datapath can still prune tokens
of the request now in flight; the two mechanisms do not substitute for
each other.
